% f18_topology_gallery variant a -- the census: six topologies drawn, each with
% how many of the suite's instances sit on it and what radix it implies.
% Data: the SUITE manifest in scripts/gen_coverage.py -- 40 instances whose
%       per-topology counts are linear 4, ring 6, mesh 4, torus 10, fat-tree 9,
%       hypercube 7 (summing to \SuiteSize{}); the same counts are the column
%       totals of tables/coverage.tex, and the N and G spans are the min/max of
%       that manifest per topology. Radix rules and ranges are tables/radix.tex
%       (linear/ring 2+local = \RadixMin; mesh/torus 4+local = \RadixLo;
%       hypercube dim+1 = \RadixMin--\RadixHc; fat-tree max(5,k+2) =
%       \RadixLo--\RadixMax), which gen_resource.py derives from the same
%       sprs_core.build_topology call the emitter used. \RadixMaxInst,
%       \RadixMaxRouters, \RouterBlowup from data/resource_macros.tex. The
%       family list T and the "one noc_system, two tables" claim are
%       sec:fabric-topo. Every drawing is schematic: the smallest member of the
%       family, not the instance -- max_bal_2d is 1,024 routers, not nine.
% Strategy: the CENSUS. One row, six shapes, and under each the three facts a
%       reader needs to place it -- how much of the suite it carries, what radix
%       it costs, and over what (N,G) span it was exercised. b puts a worked
%       route on each and argues invariance; c re-scales the panels by instance
%       count so the suite's emphasis is in the drawing; d abandons the row and
%       plots radix against per-router cost. This one is the table, drawn.
% Colour: routers and links are structural machinery (spBlue over spBlueF);
%       radix is the cost channel, so the two families whose radix GROWS with G
%       are spVerm and the four with a fixed radix are spMute; the single
%       excluded instance is a caution (spAmber); the invariance statement in
%       the footer is the confirmed result (spGreen). Shape carries the family
%       and a dashed rule under the variable-radix cells carries the cost
%       distinction, so greyscale loses nothing.
% Target: IEEE double column, 7.16in = 18.2cm. Drawn on a fixed 16.9cm canvas --
%       the standalone wrapper is varwidth=17cm, so a wider canvas is cropped in
%       the guide preview (the same constraint f04_fabric records).
\begin{tikzpicture}[
  x=1cm, y=1cm,
  rt/.style ={draw=spBlue, fill=spBlueF, circle, line width=0.35pt,
              minimum size=0.20cm, inner sep=0pt},
  sw/.style ={draw=spBlue, fill=spBlue!32, rectangle, line width=0.35pt,
              minimum width=0.28cm, minimum height=0.17cm, inner sep=0pt},
  lk/.style ={draw=spBlue!60, line width=0.4pt},
  wr/.style ={draw=spBlue!60, line width=0.4pt, densely dotted},
  ttl/.style={font=\scriptsize\bfseries, text=spInk, inner sep=1pt},
  cnt/.style={font=\scriptsize, text=spInk, inner sep=0.5pt},
  fix/.style={font=\tiny, text=spMute, inner sep=0.5pt},
  var/.style={font=\tiny, text=spVerm, inner sep=0.5pt},
  rng/.style={font=\tiny, text=spMute, inner sep=0.5pt},
]
% Fixed canvas: pins the bounding box so the preview equals the typeset size.
\path (0,0.02) rectangle (16.9,6.62);

% ---- panel grounds ---------------------------------------------------------
\begin{scope}[on background layer]
  \foreach \cx in {1.45,4.25,7.05,9.85,12.65,15.45}{
    \fill[spGrid!55, rounded corners=2.5pt] (\cx-1.35,2.24) rectangle (\cx+1.35,6.34);
  }
\end{scope}
\foreach \cx in {1.45,4.25,7.05,9.85,12.65,15.45}{
  \draw[spMute!55, line width=0.3pt] (\cx-1.20,3.60) -- (\cx+1.20,3.60);
}

% =========================== 1. LINEAR ======================================
\node[ttl] at (1.45,6.10) {linear};
\begin{scope}[shift={(1.45,4.88)}]
  \foreach \i in {0,...,4}{ \node[rt] (L\i) at (-0.88+\i*0.44,0) {}; }
  \foreach \i/\j in {0/1,1/2,2/3,3/4}{ \draw[lk] (L\i) -- (L\j); }
\end{scope}
\node[cnt] at (1.45,3.36) {\textbf{4} of \SuiteSize{} instances};
\node[fix] at (1.45,3.07) {radix $P = \RadixMin$};
\node[fix] at (1.45,2.81) {(2 + local)};
\node[rng] at (1.45,2.53) {$N$ 8--1024 \quad $G$ 2--16};

% =========================== 2. RING ========================================
\node[ttl] at (4.25,6.10) {ring};
\begin{scope}[shift={(4.25,4.88)}]
  \foreach \i in {0,...,7}{
    \node[rt] (R\i) at ({0.60*cos(90-\i*45)},{0.60*sin(90-\i*45)}) {};
  }
  \foreach \i/\j in {0/1,1/2,2/3,3/4,4/5,5/6,6/7,7/0}{ \draw[lk] (R\i) -- (R\j); }
\end{scope}
\node[cnt] at (4.25,3.36) {\textbf{6} of \SuiteSize{} instances};
\node[fix] at (4.25,3.07) {radix $P = \RadixMin$};
\node[fix] at (4.25,2.81) {(2 + local)};
\node[rng] at (4.25,2.53) {$N$ 8--2048 \quad $G$ 4--512};

% =========================== 3. MESH ========================================
\node[ttl] at (7.05,6.10) {2D mesh};
\begin{scope}[shift={(7.05,4.88)}]
  \foreach \i in {0,1,2}{ \foreach \j in {0,1,2}{
    \node[rt] (M\i\j) at (-0.54+\i*0.54,-0.54+\j*0.54) {};
  }}
  \foreach \j in {0,1,2}{ \draw[lk] (M0\j) -- (M1\j) -- (M2\j); }
  \foreach \i in {0,1,2}{ \draw[lk] (M\i0) -- (M\i1) -- (M\i2); }
\end{scope}
\node[cnt] at (7.05,3.36) {\textbf{4} of \SuiteSize{} instances};
\node[fix] at (7.05,3.07) {radix $P = \RadixLo$};
\node[fix] at (7.05,2.81) {(4 + local)};
\node[rng] at (7.05,2.53) {$N$ 16--2048 \quad $G$ 5--256};

% =========================== 4. TORUS =======================================
\node[ttl] at (9.85,6.10) {2D torus};
\begin{scope}[shift={(9.85,4.88)}]
  \foreach \i in {0,1,2}{ \foreach \j in {0,1,2}{
    \node[rt] (T\i\j) at (-0.50+\i*0.50,-0.50+\j*0.50) {};
  }}
  \foreach \j in {0,1,2}{ \draw[lk] (T0\j) -- (T1\j) -- (T2\j); }
  \foreach \i in {0,1,2}{ \draw[lk] (T\i0) -- (T\i1) -- (T\i2); }
  \foreach \j in {0,1,2}{ \draw[wr] (T0\j) to[out=150,in=30,looseness=1.1] (T2\j); }
  \foreach \i in {0,1,2}{ \draw[wr] (T\i0) to[out=240,in=120,looseness=1.1] (T\i2); }
\end{scope}
\node[cnt] at (9.85,3.36) {\textbf{10} of \SuiteSize{} instances};
\node[fix] at (9.85,3.07) {radix $P = \RadixLo$};
\node[fix] at (9.85,2.81) {(4 + local)};
\node[rng] at (9.85,2.53) {$N$ 32--8192 \quad $G$ 10--2048};

% =========================== 5. FAT-TREE ====================================
\node[ttl] at (12.65,6.10) {fat-tree};
\begin{scope}[shift={(12.65,4.88)}]
  \node[sw] (FC) at (-0.30,0.58) {};
  \node[sw] (FP0) at (-0.72,0.02) {};
  \node[sw] (FP1) at ( 0.12,0.02) {};
  \draw[lk] (FP0) -- (FC) -- (FP1);
  \foreach \i in {0,...,5}{ \node[rt] (FL\i) at (-0.98+\i*0.27,-0.58) {}; }
  \foreach \i in {0,1,2}{ \draw[lk] (FL\i) -- (FP0); }
  \foreach \i in {3,4,5}{ \draw[lk] (FL\i) -- (FP1); }
  \node[font=\tiny, text=spMute, anchor=west, inner sep=1pt] at (-0.12,0.58) {core};
  \node[font=\tiny, text=spMute, anchor=west, inner sep=1pt] at ( 0.30,0.02) {pod};
  \node[font=\tiny, text=spMute, anchor=west, inner sep=1pt] at ( 0.52,-0.58) {leaf};
\end{scope}
\node[cnt] at (12.65,3.36) {\textbf{9} of \SuiteSize{} instances};
\node[var] at (12.65,3.07) {radix $P = \RadixLo$--$\RadixMax$};
\node[var] at (12.65,2.81) {$\max(5,\,k{+}2)$};
\draw[spVerm, line width=0.4pt, densely dashed] (11.75,2.67) -- (13.55,2.67);
\node[rng] at (12.65,2.53) {$N$ 32--8192 \quad $G$ 8--4096};

% =========================== 6. HYPERCUBE ===================================
\node[ttl] at (15.45,6.10) {hypercube};
\begin{scope}[shift={(15.45,4.85)}]
  \foreach \i/\x/\y in {0/-0.62/-0.50, 1/0.20/-0.50, 2/0.20/0.32, 3/-0.62/0.32}{
    \node[rt] (H\i) at (\x,\y) {};
  }
  \foreach \i/\x/\y in {4/-0.24/-0.14, 5/0.58/-0.14, 6/0.58/0.68, 7/-0.24/0.68}{
    \node[rt] (H\i) at (\x,\y) {};
  }
  \foreach \i/\j in {0/1,1/2,2/3,3/0}{ \draw[lk] (H\i) -- (H\j); }
  \foreach \i/\j in {4/5,5/6,6/7,7/4}{ \draw[lk] (H\i) -- (H\j); }
  \foreach \i/\j in {0/4,1/5,2/6,3/7}{ \draw[wr] (H\i) -- (H\j); }
\end{scope}
\node[cnt] at (15.45,3.36) {\textbf{7} of \SuiteSize{} instances};
\node[var] at (15.45,3.07) {radix $P = \RadixMin$--$\RadixHc$};
\node[var] at (15.45,2.81) {$\dim + 1$};
\draw[spVerm, line width=0.4pt, densely dashed] (14.55,2.67) -- (16.35,2.67);
\node[rng] at (15.45,2.53) {$N$ 32--8192 \quad $G$ 4--1024};

% ---- key and footer --------------------------------------------------------
\node[font=\tiny, text=spVerm, anchor=west, inner sep=1pt] at (0.15,2.06)
  {\rule[0.06cm]{0.34cm}{0.4pt}\ \ radix grows with $G$};
\node[font=\tiny, text=spMute, anchor=west, inner sep=1pt] at (3.05,2.06)
  {radix fixed by the family};
\node[font=\tiny, text=spAmber!75!spInk, anchor=west, inner sep=1pt] at (6.00,2.06)
  {\RadixMaxInst{} ($P{=}\RadixMax$, \RadixMaxRouters{} routers) is the one
   instance the campaign excludes};

\node[anchor=north west, align=left, font=\tiny, inner sep=1pt, text=spInk,
      text width=16.6cm] at (0.15,1.80)
  {One \sv{noc\_system} module realises all six: the adjacency table maps
   (router, port) to (target router, target port) and the routing table maps
   (router, destination) to an output port, both written before
   \sv{program\_start}, so a family is a Python classmethod and no Verilog
   change. \emph{Radix is what the shape costs.} Linear and ring hold at
   $P = \RadixMin$ and mesh and torus at $P = \RadixLo$ whatever $G$ is, while
   hypercube reaches $P = \RadixHc$ at $G{=}1024$ and fat-tree $P = \RadixMax$
   at $G{=}4096$; because the crossbar is $O(P^2)$ a fat-tree router is
   \RouterBlowup$\times$ a mesh router (\RouterBitsFt{} against \RouterBitsLo{}
   storage-equivalent bits), which is why one fixed-radix estimate cannot
   price the suite.};
\node[anchor=north west, align=left, font=\tiny, inner sep=1pt, text=spGreen,
      text width=16.6cm] at (0.15,0.56)
  {\textbf{What does not vary with the shape is the answer.} Across all six
   families the target is $\Rcan(N,L)$, a function of $N$ alone: over
   \WallTBs{} settled hardware simulations no configuration produced a finite
   root differing from it.};
\end{tikzpicture}
